\documentclass[10pt,a4paper]{article}
\usepackage[margin=21mm,headheight=14pt]{geometry}
\usepackage[T1]{fontenc}
\usepackage{lmodern,microtype,amsmath,booktabs,tabularx,array}
\usepackage[dvipsnames]{xcolor}
\usepackage{graphicx,pgfplots}
\usepgfplotslibrary{groupplots}
\pgfplotsset{compat=1.18}
\usepackage{enumitem,fancyhdr,caption}
\usepackage[colorlinks=true,linkcolor=MidnightBlue,urlcolor=MidnightBlue,citecolor=MidnightBlue]{hyperref}
\definecolor{baseline}{HTML}{707D8C}
\definecolor{batcha}{HTML}{258B9C}
\definecolor{batchb}{HTML}{D6773E}
\pagestyle{fancy}\fancyhf{}
\fancyhead[L]{\small Mihal Dimo / Qwen3.8-27B on H100}
\fancyhead[R]{\small Experimental report -- 1 October 2026}
\fancyfoot[C]{\small\thepage}
\setlength{\parindent}{0pt}\setlength{\parskip}{5pt}
\setlength{\tabcolsep}{5pt}
\renewcommand{\arraystretch}{1.15}
\captionsetup{font=small,labelfont=bf}
\setlist[itemize]{leftmargin=*,itemsep=2pt,topsep=2pt}
\newcommand{\code}[1]{\texttt{\small #1}}
% Generated from raw results; do not hand edit.
\newcommand{\BaseSingle}{74.31}
\newcommand{\BaseSingleTTFT}{62.89 / 63.91}
\newcommand{\BaseSingleTTFTMedian}{62.89}
\newcommand{\BaseSingleTPOT}{13.07 / 13.07}
\newcommand{\BaseSingleTPOTMedian}{13.07}
\newcommand{\BasePrefill}{11,919.10}
\newcommand{\BasePrefillTTFT}{10,978.87 / 11,010.75}
\newcommand{\BasePrefillTTFTMedian}{10,978.87}
\newcommand{\BaseDecode}{74.37}
\newcommand{\BaseDecodeTTFT}{53,413.02 / 53,425.81}
\newcommand{\BaseDecodeTTFTMedian}{53,413.02}
\newcommand{\BaseDecodeTPOT}{13.06 / 13.06}
\newcommand{\BaseDecodeTPOTMedian}{13.06}
\newcommand{\BaseMixed}{71.70}
\newcommand{\BaseMixedTTFT}{55,448.50 / 55,475.93}
\newcommand{\BaseMixedTTFTMedian}{55,448.50}
\newcommand{\BaseMixedTPOT}{13.09 / 13.10}
\newcommand{\BaseMixedTPOTMedian}{13.09}
\newcommand{\BaseSpread}{0.05}
\newcommand{\BaseMemory}{33.36}
\newcommand{\BasePower}{680.5--685.1}
\newcommand{\SmallSingle}{72.83}
\newcommand{\SmallSingleTTFT}{98.54 / 100.09}
\newcommand{\SmallSingleTTFTMedian}{98.54}
\newcommand{\SmallSingleTPOT}{13.06 / 13.06}
\newcommand{\SmallSingleTPOTMedian}{13.06}
\newcommand{\SmallPrefill}{17,683.68}
\newcommand{\SmallPrefillTTFT}{7,329.12 / 7,536.07}
\newcommand{\SmallPrefillTTFTMedian}{7,329.12}
\newcommand{\SmallDecode}{1,647.90}
\newcommand{\SmallDecodeTTFT}{335.37 / 339.57}
\newcommand{\SmallDecodeTTFTMedian}{335.37}
\newcommand{\SmallDecodeTPOT}{16.92 / 16.94}
\newcommand{\SmallDecodeTPOTMedian}{16.92}
\newcommand{\SmallMixed}{996.70}
\newcommand{\SmallMixedTTFT}{1,519.94 / 1,899.09}
\newcommand{\SmallMixedTTFTMedian}{1,519.94}
\newcommand{\SmallMixedTPOT}{20.26 / 26.28}
\newcommand{\SmallMixedTPOTMedian}{20.26}
\newcommand{\SmallSpread}{0.09}
\newcommand{\SmallMemory}{48.81}
\newcommand{\SmallPower}{678.0--681.1}
\newcommand{\LargeSingle}{72.49}
\newcommand{\LargeSingleTTFT}{98.00 / 98.75}
\newcommand{\LargeSingleTTFTMedian}{98.00}
\newcommand{\LargeSingleTPOT}{13.13 / 13.14}
\newcommand{\LargeSingleTPOTMedian}{13.13}
\newcommand{\LargePrefill}{18,092.46}
\newcommand{\LargePrefillTTFT}{7,065.13 / 7,997.89}
\newcommand{\LargePrefillTTFTMedian}{7,065.13}
\newcommand{\LargeDecode}{1,645.66}
\newcommand{\LargeDecodeTTFT}{333.70 / 335.75}
\newcommand{\LargeDecodeTTFTMedian}{333.70}
\newcommand{\LargeDecodeTPOT}{16.96 / 16.99}
\newcommand{\LargeDecodeTPOTMedian}{16.96}
\newcommand{\LargeMixed}{1,003.64}
\newcommand{\LargeMixedTTFT}{1,779.62 / 1,856.21}
\newcommand{\LargeMixedTTFTMedian}{1,779.62}
\newcommand{\LargeMixedTPOT}{18.11 / 21.45}
\newcommand{\LargeMixedTPOTMedian}{18.11}
\newcommand{\LargeSpread}{0.56}
\newcommand{\LargeMemory}{50.97}
\newcommand{\LargePower}{681.3--685.0}
\newcommand{\DecodeSpeedup}{22.13}
\newcommand{\MixedSpeedup}{14.00}
\newcommand{\PrefillRatio}{1.52}
\newcommand{\PrefillGain}{51.79}
\newcommand{\SingleRatio}{0.98}
\newcommand{\SingleLoss}{2.45}
\newcommand{\BudgetPrefillGain}{2.31}
\newcommand{\BudgetMixedGain}{0.70}
\newcommand{\BudgetDecodeChange}{-0.14}

\title{\vspace{-12mm}\textbf{From Single-Request Inference to Concurrent Serving}\\[3pt]
\large Throughput--Latency Tradeoffs for Qwen3.8-27B on an H100}
\author{Mihal Dimo\\\small\href{https://github.com/mhdimo/qwen38-h100-lab}{github.com/mhdimo/qwen38-h100-lab}}
\date{1 October 2026}

\begin{document}
\maketitle\thispagestyle{fancy}
\vspace{-6mm}
\begin{abstract}
We evaluate an isolated concurrent-inference implementation for Qwen3.8-27B-FP8
on one NVIDIA H100 80GB GPU. Inspection of the existing deployment reveals
a one-sequence engine limit, a 2 GiB cache budget, and an application lock that
rejects overlapping inference. Our implementation uses vLLM AsyncLLM with a
64-sequence limit, 16 GiB cache, and CUDA-graph capture sizes for concurrent
decoding. With 32 concurrent synthetic requests, the 16,384-token scheduling
budget increases decode-heavy aggregate output throughput from
\BaseDecode{} to \LargeDecode{} tokens/s (\DecodeSpeedup{}$\times$), and mixed
throughput from \BaseMixed{} to \LargeMixed{} tokens/s (\MixedSpeedup{}$\times$).
Prefill-heavy input throughput rises by \PrefillGain{}\%.
These gains trade higher memory use and per-request token latency for much less
queueing. Single-request first-token latency regresses. The comparison is between
matched engine configurations under an asynchronous harness, not an HTTP
production benchmark or a novel GPU-kernel implementation. Three configurations
completed; later experiments were cancelled.
\end{abstract}

\section{System question and contribution}
The objective is maximum aggregate throughput for many concurrent users, rather
than minimum latency for one request. vLLM exposes batching and memory-management
mechanisms for this purpose~\cite{pagedattention}. However, the existing
application's admission policy prevents concurrent requests from reaching the
engine. Optimizing a GPU kernel alone cannot remove that constraint.

The contribution of this report is a reproducible local benchmark implementation,
a measured comparison of engine policies, and a documented integration boundary.
Across measured configurations we preserve the checkpoint, numerical formats,
runtime and context limit. We do not claim a new attention algorithm, custom CUDA kernel, or universal
fastest configuration.

\begin{table}[h]
\centering\small
\caption{Deployed worker versus the proposed implementation.}
\begin{tabularx}{\linewidth}{@{}l>{\raggedright\arraybackslash}X>{\raggedright\arraybackslash}X@{}}
\toprule
Property & Current deployed worker & Proposed isolated implementation \\
\midrule
Request path & Synchronous \code{LLM.generate}; application busy lock & Bounded asynchronous request pool feeding \code{AsyncLLM} \\
Sequence capacity & 1 & 64 \\
Fixed cache budget & 2 GiB & 16 GiB \\
Prefill token budget & 2,048 in reproduced baseline & 8,192 or 16,384 \\
Compilation / graphs & Enabled; prefill-oriented captures through 2,048 & Enabled; captures through batch size 64 \\
Deployment status & Existing service & Separate lab; not integrated into service \\
\bottomrule
\end{tabularx}
\end{table}

\newpage
\section{Experimental method}
\subsection{Hardware, model and software}
All inference ran on one isolated NVIDIA H100 80GB HBM3 with an observed
700 W power limit and driver 610.57.04. The software stack was Python 3.14.6,
vLLM 0.26.0, PyTorch 2.11.0, Triton 3.6.0, Transformers 5.16.1 and FlashInfer
0.6.14. The cached checkpoint was \code{Qwen/Qwen3.8-27B-FP8}.\\
Pinned revision: \code{017b9c7af6b5689d5dd426a76e0bc077eb5ca20a}.
Qwen3.8-27B combines Gated DeltaNet and full-attention layers in the Qwen3.5
architecture~\cite{qwen}. The native FP8 checkpoint and automatic BF16 cache
format were unchanged; maximum context length was 10,240 tokens.

Resolved logs identify FP8 DeepGEMM, FlashAttention 3, FlashInfer GDN prefill,
asynchronous engine scheduling, and norm/activation quantization fusions.
This is evidence of selected implementations, not a measurement of their
individual time shares.

\subsection{Baseline construction and fairness}
Every measured configuration uses the same direct AsyncLLM client. The baseline
reproduces the deployed engine's one-sequence limit, 2 GiB cache and graph policy.
At client concurrency 32, its engine queues requests. The actual deployed wrapper
instead raises a busy error when another inference is running. Therefore the
baseline is an \emph{engine-capacity proxy}, not a measurement of accepted HTTP
traffic, network latency or production error rate.

The first candidate changes sequence capacity, cache capacity and graph capture
policy together. Its gain must be attributed to that combined configuration.
The 8,192-versus-16,384 comparison changes only the scheduled-token budget.
Graph mode is \code{FULL\_AND\_PIECEWISE}; candidate capture sizes are
\{1,2,4,8,16,24,32,48,64\}. The baseline uses
\{1,2,128,160,256,512,1024,2048\}.

\begin{table}[h]
\centering\small
\caption{Fixed-length synthetic workloads. Each configuration has two repetitions.}
\begin{tabular}{@{}lrrrr@{}}
\toprule
Workload & Input tokens & Output tokens & Concurrency & Requests / repetition \\
\midrule
Single-request decode & 128 & 128 & 1 & 8 \\
Prefill-heavy & 4,096 & 1 & 32 & 64 \\
Decode-heavy & 128 & 128 & 32 & 64 \\
Mixed & 1,024 & 128 & 32 & 64 \\
\bottomrule
\end{tabular}
\end{table}

Prompts are seeded, unique token-ID sequences, shared across candidates.
Fixed output length is enforced with \code{ignore\_eos}, \code{min\_tokens} and
\code{max\_tokens}. Warmup uses separate prompts and the supported concurrency
shape, followed by prefix-cache reset. All reported measured requests have zero
cached prompt tokens. These inputs exercise tensor shapes, not representative
language or vision semantics. Each complete configuration contributes 400 measured
requests; the three complete configurations contribute 1,200.

\subsection{Metrics and repetition handling}
Let $T$ be wall-clock duration from the first client dispatch until all requests
complete. For request $i$, let $s_i$, $f_i$ and $e_i$ denote submission, first
observed output, and completion; $P_i$ and $O_i$ are token counts. We report
\[
R_{\rm out}=\frac{\sum_i O_i}{T},\quad
R_{\rm in}=\frac{\sum_i P_i}{T},\quad
\mathrm{TTFT}_i=f_i-s_i,\quad
\mathrm{TPOT}_i=\frac{e_i-f_i}{O_i-1}.
\]
TPOT is undefined for one-token outputs. It is a client-observed request average,
not an individual inter-token latency distribution. Repetition throughput pools
token totals and wall times; p50/p95 pool per-request samples. Two repetitions
support descriptive repeatability, not confidence intervals or significance tests.

\newpage
\section{Measured throughput}
\begin{table}[h]
\centering\small
\caption{Completed runs only. Output throughput unless the row says input.}
\begin{tabular}{@{}lrrrr@{}}
\toprule
Workload / metric & Baseline & 8,192 budget & 16,384 budget & 16,384 / baseline \\
\midrule
Single-request output tok/s & \BaseSingle & \SmallSingle & \LargeSingle & \SingleRatio{}$\times$ \\
Prefill input tok/s & \BasePrefill & \SmallPrefill & \LargePrefill & \PrefillRatio{}$\times$ \\
Decode-heavy output tok/s & \BaseDecode & \SmallDecode & \LargeDecode & \DecodeSpeedup{}$\times$ \\
Mixed output tok/s & \BaseMixed & \SmallMixed & \LargeMixed & \MixedSpeedup{}$\times$ \\
\bottomrule
\end{tabular}
\end{table}

\begin{figure}[h]
\centering
\begin{tikzpicture}
\begin{groupplot}[group style={group size=2 by 1,horizontal sep=16mm},
 width=.44\linewidth,height=58mm,
 ybar,/pgf/bar width=11pt,enlarge x limits=.3,
 tick label style={font=\footnotesize},label style={font=\small},
 axis line style={gray!50},ymajorgrids,grid style={gray!15},
 symbolic x coords={Base,8k,16k},xtick=data,ymin=0,
 legend style={font=\footnotesize,draw=none},scaled y ticks=false,
 /pgf/number format/1000 sep={}]
\nextgroupplot[ylabel={Output tokens/s},title={Concurrent generation},ymax=2000,
 legend to name=throughputlegend,legend columns=2]
\addplot[fill=batcha,draw=none] table[x=config,y=decode,col sep=comma]{throughput.csv};
\addplot[fill=batchb,draw=none] table[x=config,y=mixed,col sep=comma]{throughput.csv};
\legend{Decode-heavy,Mixed}
\nextgroupplot[ylabel={Input tokens/s},title={Prefill-heavy},ymax=22000]
\addplot[fill=baseline,draw=none] table[x=config,y=prefill,col sep=comma]{throughput.csv};
\end{groupplot}
\end{tikzpicture}
\vspace{1mm}\ref{throughputlegend}
\caption{Matched concurrency-32 measurements. Prefill-heavy requests produce one
output token; their rate includes scheduling and first-token generation, and is
not isolated GPU prefill speed.}
\end{figure}

\subsection{Batching dominates the first optimization step}
The baseline can execute only one sequence at a time. Enlarging engine capacity
allows concurrent requests to share execution steps, raising decode-heavy output
throughput by \DecodeSpeedup{}$\times$. Mixed throughput improves by
\MixedSpeedup{}$\times$. The prefill-heavy gain is smaller
(\PrefillGain{}\%), consistent with a workload that already presents substantial
token-level parallelism. This interpretation follows the configuration contrast;
kernel-level causal attribution would require the cancelled profiles.

\subsection{The larger token budget gives a modest incremental gain}
Moving from 8,192 to 16,384 increases prefill input throughput by
\BudgetPrefillGain{}\% and mixed output throughput by \BudgetMixedGain{}\%,
while decode-heavy throughput changes by \BudgetDecodeChange{}\%.
These differences are small compared with enabling concurrency. No single
budget is uniformly superior: the 8,192 setting has lower median mixed TTFT and
lower observed memory use, while 16,384 has higher aggregate prefill throughput.

For the two measured repetitions, the largest relative output-rate spread across
the four cases is \SmallSpread{}\% for 8,192 and \LargeSpread{}\% for 16,384,
calculated as $(\max-\min)/\operatorname{mean}$. We do not interpret these
two-run ranges as uncertainty bounds on production workloads.

\newpage
\section{Latency and memory tradeoffs}
\begin{table}[h]
\centering\small
\caption{Client-observed latency: p50 / p95. Values are milliseconds.}
\begin{tabular}{@{}llrrr@{}}
\toprule
Workload & Metric & Baseline & 8,192 budget & 16,384 budget \\
\midrule
Single request & TTFT & \BaseSingleTTFT & \SmallSingleTTFT & \LargeSingleTTFT \\
Single request & TPOT & \BaseSingleTPOT & \SmallSingleTPOT & \LargeSingleTPOT \\
Prefill-heavy & TTFT & \BasePrefillTTFT & \SmallPrefillTTFT & \LargePrefillTTFT \\
Decode-heavy & TTFT & \BaseDecodeTTFT & \SmallDecodeTTFT & \LargeDecodeTTFT \\
Decode-heavy & TPOT & \BaseDecodeTPOT & \SmallDecodeTPOT & \LargeDecodeTPOT \\
Mixed & TTFT & \BaseMixedTTFT & \SmallMixedTTFT & \LargeMixedTTFT \\
Mixed & TPOT & \BaseMixedTPOT & \SmallMixedTPOT & \LargeMixedTPOT \\
\bottomrule
\end{tabular}
\end{table}

\textbf{Throughput versus token latency.} Decode-heavy TPOT p50 increases from
\BaseDecodeTPOTMedian{} to \LargeDecodeTPOTMedian{} ms. Each concurrent request
therefore receives tokens more slowly than a request already executing alone.
Aggregate output rate nevertheless increases sharply because many requests make
progress together. Baseline median decode TTFT is approximately 53.4 s because
requests queue; the 16,384 candidate reduces it to approximately 0.334 s.
These finite closed-loop experiments include startup/drain effects and are not
open-loop service-level capacity tests.

\textbf{A low-load regression.} At concurrency one, median TTFT rises from
\BaseSingleTTFTMedian{} to \LargeSingleTTFTMedian{} ms and output throughput
declines by \SingleLoss{}\%. The graph capture policies differ, so this result
does not isolate its cause. A throughput-oriented policy should not be presented
as a universal latency improvement.

\begin{table}[h]
\centering\small
\caption{Resource cost. GPU memory is the maximum of end-of-case snapshots,
not a continuously sampled peak. All runs use the same 700 W power limit.}
\begin{tabular}{@{}lrrr@{}}
\toprule
Resource & Baseline & 8,192 budget & 16,384 budget \\
\midrule
Configured cache (GiB) & 2 & 16 & 16 \\
Observed GPU memory (GiB) & \BaseMemory & \SmallMemory & \LargeMemory \\
Prefill end-of-case power (W), range & \BasePower & \SmallPower & \LargePower \\
\bottomrule
\end{tabular}
\end{table}

The larger cache and working set reduce headroom for other model workers. The
16,384 candidate reaches \LargeMemory{} GiB in sampled GPU memory, compared with
\BaseMemory{} GiB for the baseline. Co-resident model capacity must be re-evaluated
before adoption. Prefill samples approach the power limit; clocks were not locked
and varied with load. Power snapshots do not establish energy per request.

\subsection{Compilation and CUDA graphs}
Kernel compilation and CUDA graphs solve complementary problems: compilation
optimizes generated GPU work, while graph replay reduces repeated launch
overhead~\cite{cuda,vllmgraphs}. Both were already enabled in the baseline.
Cached compiled artifacts reduced observed compilation initialization from
58.45 s on an earlier startup attempt to 7.63 s in the successful baseline
restart. This is a startup observation, not a steady-state throughput speedup
or a controlled startup benchmark. All serving measurements exclude engine
initialization and workload warmup. A graphs-off ablation was not performed.

\newpage
\section{Validity, quality and deployment boundary}
\subsection{What the results establish}
The completed configurations passed three natural-language smoke checks
(capital of France, $7\times8$, opposite of hot), completed every measured request
at the required output length, and reported zero cached prompt tokens. The
complete comparison uses unchanged checkpoint bytes and package versions.
The benchmark's CPU tests cover wall-time throughput, cumulative streaming,
single-token TPOT, concurrency bounds, incomplete output, timeout evidence and
process cleanup; 16 tests passed locally.

These checks do not establish numerical equivalence or production model quality.
Even the smoke outputs differ in capitalization (``cold'' versus ``Cold'').
Synthetic token sequences are not meaningful quality examples. No image/video,
structured-output, long-conversation or task-accuracy evaluation was completed.
Backend arithmetic, batch shape and graph padding can affect generated tokens;
the report does not claim identical outputs.

\subsection{Threats to validity}
\begin{itemize}
\item \textbf{Traffic model.} Only fixed lengths and concurrency 1/32 were measured
fully. There is no variable-length arrival distribution, sustained overload,
admission-rate sweep or production network path.
\item \textbf{Attribution.} The main comparison jointly changes cache size,
sequence capacity and graph policy. It cannot isolate batching, cache or graphs
as the sole source of speedup.
\item \textbf{Small repetition count.} Two repetitions per case provide limited
evidence about stability. Trials were not randomized or interleaved, and the
successful baseline rerun occurred after the two larger-budget candidates.
\item \textbf{Timing semantics.} TTFT includes queueing; TPOT is based on streamed
observations. Grouped first chunks can obscure true token-generation timing.
GPU kernel timings and a kernel-level bottleneck breakdown are unavailable.
\item \textbf{Resources.} Memory and power are point samples. The model manager
remained active, but no other GPU compute process was present at experiment
startup. No concurrent co-residency test or energy measurement was performed.
\end{itemize}

\subsection{Final machine state}
No remaining live lab processes were found. A subsequent GPU query returned
no compute processes, while the production inference service remained active
and untouched. The experiment folder, logs and raw results were copied to a
local archive. No production configuration or model checkpoint was changed.

\section{Recommendation}
For the measured concurrency-32 workloads, retain the 16,384-token candidate as
the throughput-oriented reference and 8,192 as the lower-memory, lower-median-
mixed-TTFT alternative. Neither is a validated production default. To expose
the measured engine capacity to users, a future integration must replace the
single-request busy lock with asynchronous admission and bounded backpressure,
preserve cancellation and structured/multimodal behavior, and update memory
accounting for co-resident models.

The next kernel step is profiling after warmup, followed by one backend or kernel
ablation with a numerical reference and matched end-to-end tests~\cite{cuda}.
Compilation and graph capture should remain part of that procedure. Current
evidence supports a substantial engine-level throughput gain; it does not yet
support a claim of custom-kernel acceleration or saturation of the H100.

\newpage
\appendix
\section{Reproduction and artifact map}
\textbf{Archive root:} this repository.\\
\textbf{GPU workspace:} a temporary lab directory on the H100 host, removed
after hash-verified transfer. Inference is stopped. Generating this paper and analyzing its data require no GPU.

\begin{table}[h]
\centering\small
\caption{Source evidence for the reported configurations.}
\begin{tabularx}{\linewidth}{@{}lX@{}}
\toprule
Configuration & Raw JSON under \code{results/} \\
\midrule
Baseline & \nolinkurl{20261001T113011.178231Z/baseline.json} \\
8,192-token candidate & \nolinkurl{20261001T102025.566538Z/batch64_t8192.json} \\
16,384-token candidate & \nolinkurl{20261001T102025.566538Z/batch64_t16384.json} \\
Interrupted candidate & \nolinkurl{20261001T113011.178231Z/batch64_t2048.json} \\
\bottomrule
\end{tabularx}
\end{table}

Each result stores per-request token counts and timings, output-token hashes,
workload definitions, package versions, hardware snapshots, model path, resolved
engine configuration and source hashes. Initialization logs preserve selected
backends, compile/cache messages and GPU-cache sizing. Two initial attempts
failed because \code{ninja} was absent from the launch PATH; the installed Python
binary directory was added without changing dependencies. The broad cancellation
marker also tagged these earlier incomplete JSON files as cancelled; their
nonzero exit codes and original failure logs identify them as startup failures.
They are excluded throughout.

\textbf{Build this paper locally:}
\begin{verbatim}
python3 paper/build_data.py
cd paper
tectonic paper.tex
\end{verbatim}

\textbf{Benchmark entry points, retained for a later authorized GPU run:}
\begin{verbatim}
python run_suite.py --configs baseline batch64_t8192 batch64_t16384
python run_suite.py --configs batch128_t16384 \
  --workloads workloads_scale.json
python run_suite.py --configs batch64_t16384 \
  --workloads workloads_profile.json --repetitions 1 --profile
\end{verbatim}
On the H100 host, \code{python} above means the pinned Python 3.14.6 installation.
\code{run\_suite.py} acquires a lab lock, checks GPU idleness and serializes runs.
Do not execute the retained queue script merely to rebuild this paper.

\begin{thebibliography}{9}\small
\bibitem{pagedattention} W. Kwon et al.
\emph{Efficient Memory Management for Large Language Model Serving with
PagedAttention}. SOSP, 2023. \url{https://arxiv.org/abs/2309.06180}.
\bibitem{qwen} Qwen Team. \emph{Qwen3.8-27B model card}, 2026.
\url{https://huggingface.co/Qwen/Qwen3.8-27B}.
\bibitem{vllmopt} vLLM Project. \emph{Optimization and Tuning}.
\url{https://docs.vllm.ai/en/latest/configuration/optimization/}.
\bibitem{vllmgraphs} vLLM Project. \emph{CUDA Graphs} and
\emph{torch.compile integration}.
\url{https://docs.vllm.ai/en/latest/design/cuda_graphs/};
\url{https://docs.vllm.ai/en/latest/design/torch_compile/}.
\bibitem{cuda} NVIDIA. \emph{CUDA C++ Best Practices Guide}.
\url{https://docs.nvidia.com/cuda/cuda-c-best-practices-guide/}.
\end{thebibliography}
\small Online references consulted on 1 October 2026. Installed runtime logs,
rather than newer online defaults, define the configurations measured here.
\end{document}
